\documentclass[11pt,a4paper]{article}
\usepackage[margin=2.6cm]{geometry}
\usepackage{amsmath,amssymb,amsthm}
\usepackage{booktabs}
\usepackage[hidelinks]{hyperref}
\newtheorem{theorem}{Theorem}
\newtheorem{proposition}[theorem]{Proposition}
\title{The direction spectrum of no-three-in-line solutions:\\ a line model derives its shape, not its scale}
\author{Aleksei Kudriashov (Alex Komang)\\ \small Nusa Dua Studio, Nusa Dua, Bali, Indonesia\\ \small ORCID: 0009-0006-8885-5338}
\date{3 September 2026 --- version 1.0}
\begin{document}\maketitle
\begin{abstract}
A \emph{solution} of the no-three-in-line problem is a set of $2n$ points of the $n\times n$ grid with no three
collinear.  For a primitive direction $v$ let $c_v(n)$ be the mean number of point pairs of a solution whose
difference is a positive multiple of $v$, divided by $n$ (averaged over the images of $v$ under the symmetries of
the square).  Measured on Flammenkamp's database of all known solutions, the short directions are nearly constant in
$n$ ($c_{(1,1)}=0.731$, $c_{(1,2)}=0.564$) and depleted against a null model with fixed row and column sums
(ratio $0.58$ and $0.75$), while long directions are slightly enriched and grow with $n$.  We show that the
\emph{shape} of this spectrum --- the ratios between directions and their order --- follows from a single rule with
no fitted parameter: distribute $2n$ points over the lines of one direction with at most two points per line, the
weight of a line of $L$ cells carrying $x$ points being $\binom{L}{x}$.  The expected number of doubly occupied lines
is given exactly by a generating function; one global renormalisation to the exact number of non-axial pairs of a
solution, $2n^2-3n$, fixes the scale.  The model reproduces fifteen measured constants within $12\%$ (model$/$data $0.89$--$1.05$ at $n=20$,
$0.88$--$1.02$ at $n\approx30$) and their order at matched $n$ up to one adjacent pair; seven of them were predicted before being measured by a second agent who had not seen the
model.  What the model does \emph{not} do is stated with equal care: the scale is imposed, not derived --- the
model's own total falls $17\%$ short of $2n^2-3n$, and the shortfall it redistributes is exactly the deficit seen in the
measured directions; the growth of long directions with $n$ is a property of the solutions that the model only
partly tracks; and no prior-art search beyond the classical sources has yet been run.
\end{abstract}

\section{Definitions and data}
Let $[n]=\{0,\dots,n-1\}$.  A set $S\subseteq[n]^2$ is \emph{admissible} if no three of its points are collinear
and a \emph{solution} if moreover $|S|=2n$, the maximum possible since every row carries at most two points.  For
$v=(a,b)$ with $\gcd(a,b)=1$ and $(a,b)\ne(0,0)$, a \emph{pair of direction $v$} in $S$ is an unordered pair
$\{p,q\}\subseteq S$ with $q-p=kv$ for some $k\ge1$ (so every pair has exactly one primitive direction up to sign).
For a solution, every line carries at most two points, hence the number of pairs of direction $v$ equals the number of
lines of direction $v$ carrying exactly two points.  Define
\[c_v(n)=\frac{1}{n}\,\mathbb{E}\bigl[\#\{\text{pairs of direction }v\}\bigr],\]
the expectation over the uniform distribution on the solutions of size $n$ in the sample at hand, averaged over the
class $\{(a,b),(b,a),(a,-b),(b,-a)\}$ of $v$ under the symmetries of the square.  Since every row and every column of a
solution carries exactly two points, $c_{(1,0)}=c_{(0,1)}=1$ identically, and the total over all non-axial classes
is fixed:
\begin{equation}\label{eq:total}
\sum_{v\ \text{non-axial}} n\,c_v(n)\cdot|\text{class}(v)| = \binom{2n}{2}-2n = 2n^2-3n .
\end{equation}

\emph{Data.}  We use Achim Flammenkamp's database of all known solutions (state of 31 August 2026).  It is complete
for $n\le20$ (for $n=19$: $32\,577$ solutions; $n=20$: $118\,057$) and contains only solutions with a non-trivial
symmetry for $21\le n\le57$.  The measurements were made by a second agent with code written independently of the
model below; the values quoted as ``mixed $n$'' average over all $391\,812$ solutions with weight one.

\begin{center}\begin{tabular}{lrrrr}\toprule
class & mixed $n$ & $n=19$--$20$ & $n=29$--$31$ & $n=32$--$57$\\\midrule
$(1,1)$ & 0.7306 & 0.7308 & 0.7327 & 0.7325\\
$(1,2)$ & 0.5633 & 0.5638 & 0.5635 & 0.5548\\
$(1,3)$ & 0.4545 & 0.4545 & 0.4586 & 0.4526\\
$(2,3)$ & 0.4095 & 0.4097 & 0.4121 & 0.4070\\
$(1,4)$ & 0.3744 & 0.3741 & 0.3807 & 0.3793\\
$(1,5)$ & 0.3073 & 0.3074 & 0.3228 & 0.3261\\
$(2,5)$ & 0.2846 & 0.2842 & 0.3023 & 0.3049\\
$(3,4)$ & 0.3043 & 0.3039 & 0.3175 & 0.3186\\
$(1,6)$ & 0.2618 & 0.2492 & 0.2687 & 0.2800\\
$(1,7)$ & 0.2233 & 0.2071 & 0.2340 & 0.2461\\
$(2,7)$ & 0.2081 & 0.1889 & 0.2208 & 0.2335\\
$(3,5)$ & 0.2701 & 0.2609 & 0.2758 & 0.2819\\
$(4,5)$ & 0.2390 & 0.2243 & 0.2486 & 0.2575\\
$(1,8)$ & 0.1859 & 0.1692 & 0.1954 & 0.2160\\
$(3,7)$ & 0.1961 & 0.1791 & 0.2069 & 0.2209\\
\end{tabular}\end{center}
Short directions are constant in $n$ to $1\%$; long directions grow with $n$ (by $23\%$ for $(3,7)$ and $28\%$ for
$(1,8)$ between $n\approx20$ and $n\ge32$); the sizes $21\le n\le28$ ($54\,604$ solutions) are omitted from the table.  A ``constant of the database'' for a long direction is therefore an average over a mixture
of sizes, and the model must be compared with the data \emph{at the same $n$}.  Symmetry does not explain the growth:
at $n=19$ and $20$, where the database is complete, the spectrum of the solutions with a half-turn symmetry
($592$ and $675$ of them) differs from that of all solutions by at most $4\%$, and for the long directions in the
\emph{downward} direction ($c_{(1,7)}$: $0.1986$ against $0.2067$ at $n=19$).

\emph{The null model.}  Under the uniform distribution on all configurations with exactly two points in every row
and column (the correct null model: a solution is such a configuration), the expected number of pairs of direction
$v$ is the number of cell pairs of direction $v$ times the probability $2(2n-3)/(n(n-1)^2)$ that a given pair of cells
in distinct rows and columns is occupied (the group $S_n\times S_n$ acts transitively on such pairs).  At $n=20$ this
gives $c^{\rm null}_{(1,1)}=1.27$, $c^{\rm null}_{(1,2)}=0.75$, $\dots$, so the solutions are depleted in the short
directions ($0.58$, $0.75$, $0.89$, $0.94$ for $(1,1)$, $(1,2)$, $(1,3)$, $(2,3)$) and slightly enriched in the long ones
($1.01$--$1.11$).  This is the pattern to be explained.

\section{The line model}
Fix a direction $v$.  The cells of $[n]^2$ are partitioned into the lines of direction $v$ (a cell without a
neighbour in direction $v$ is a line of length $1$); let the lines have lengths $L_1,\dots,L_r$.

\begin{theorem}\label{thm:gf}
Let $2\le m\le\sum_i\min(L_i,2)$ and let $S$ be uniformly distributed on the $m$-subsets of the cells that contain at
most two cells of every line.  Then the number of such subsets is $Z_m=[t^m]\prod_{i=1}^{r}\bigl(1+L_it+\binom{L_i}{2}t^2\bigr)$
and
\[\mathbb{E}\bigl[\#\{i:\ |S\cap\ell_i|=2\}\bigr]=\frac{1}{Z_m}\sum_{i=1}^{r}\binom{L_i}{2}\,[t^{m-2}]\prod_{j\ne i}\bigl(1+L_jt+\binom{L_j}{2}t^2\bigr).\]
\end{theorem}
\begin{proof}
A subset is admissible for $v$ exactly when it takes $x_i\in\{0,1,2\}$ cells on line $i$; there are $\binom{L_i}{x_i}$
ways to choose them, independently across lines, so the admissible $m$-subsets are counted by the coefficient of $t^m$
in the product.  The expectation is $\sum_iP(x_i=2)$, and the admissible $m$-subsets with $x_i=2$ are counted by
$\binom{L_i}{2}$ times the coefficient of $t^{m-2}$ in the product over $j\ne i$.
\end{proof}

This is the canonical ensemble of hard particles in $r$ independent boxes of capacity two with degeneracies
$1,L_i,\binom{L_i}{2}$ (equivalently, a uniform placement conditioned on ``at most two per box''); the identity is
elementary.  What is not elementary is that it describes the solutions of the no-three-in-line problem.

\emph{The model.}  Put $m=2n$ and let $E_v$ be the expectation of Theorem~\ref{thm:gf} for the lines of direction $v$.
The model does not know that the pairs of a solution are counted by \eqref{eq:total}: summed over all non-axial
primitive directions, $\sum_vE_v=615.3$ at $n=20$ against the true $740$, a shortfall of $17\%$.  We therefore
impose \eqref{eq:total} by one global factor, $Z_n=(2n^2-3n)/\sum_vE_v$ ($Z_{20}=1.203$, $Z_{30}=1.163$,
$Z_{40}=1.139$), and set $c^{\rm model}_v(n)=Z_nE_v/n$, averaged over the class of $v$.  The factor is not fitted to
the measured spectrum --- it is the same for every direction --- but it is not derived either.  Everything the model
predicts is therefore a statement about \emph{ratios} between directions.

\section{Results}
\begin{center}\begin{tabular}{lrrrrrr}\toprule
class & model $20$ & data $19$--$20$ & ratio & model $30$ & data $29$--$31$ & ratio\\\midrule
$(1,1)$ & 0.7637 & 0.7308 & 1.045 & 0.7441 & 0.7327 & 1.016\\
$(1,2)$ & 0.5629 & 0.5638 & 0.998 & 0.5541 & 0.5635 & 0.983\\
$(1,3)$ & 0.4327 & 0.4545 & 0.952 & 0.4316 & 0.4586 & 0.941\\
$(2,3)$ & 0.3766 & 0.4097 & 0.919 & 0.3792 & 0.4121 & 0.920\\
$(1,4)$ & 0.3428 & 0.3741 & 0.916 & 0.3492 & 0.3807 & 0.917\\
$(1,5)$ & 0.2781 & 0.3074 & 0.905 & 0.2890 & 0.3228 & 0.895\\
$(2,5)$ & 0.2544 & 0.2842 & 0.895 & 0.2665 & 0.3023 & 0.882\\
$(3,4)$ & 0.2737 & 0.3039 & 0.901 & 0.2839 & 0.3175 & 0.894\\
$(1,6)$ & 0.2322 & 0.2492 & 0.932 & 0.2440 & 0.2687 & 0.908\\
$(1,7)$ & 0.1931 & 0.2071 & 0.932 & 0.2105 & 0.2340 & 0.900\\
$(2,7)$ & 0.1800 & 0.1889 & 0.953 & 0.1978 & 0.2208 & 0.896\\
$(3,5)$ & 0.2309 & 0.2609 & 0.885 & 0.2443 & 0.2758 & 0.886\\
$(4,5)$ & 0.2076 & 0.2243 & 0.926 & 0.2223 & 0.2486 & 0.894\\
$(1,8)$ & 0.1657 & 0.1692 & 0.979 & 0.1826 & 0.1954 & 0.934\\
$(3,7)$ & 0.1669 & 0.1791 & 0.932 & 0.1852 & 0.2069 & 0.895\\
\bottomrule\end{tabular}\end{center}
At $n=20$ the model's order of the fifteen classes coincides with the data's except for the adjacent pair $(1,6)$, $(3,5)$
(which differ by $3\%$ in the data); the model is below the data in every class but $(1,1)$, by $5$--$12\%$, and at
$n\approx30$ by $6$--$12\%$ in every class but $(1,1)$ and $(1,2)$.

\emph{Blind test.}  The first eight classes were known to us when the model was built.  The remaining seven were
\emph{not}: their model values were written into a dated ledger (3 September 2026, 16:05 WITA) before a second agent,
who had not seen the model or the ledger, measured them on the database with independent code; all seven fell within
$2$--$10\%$ of the prediction (model$/$data at mixed $n$: $0.93$, $0.95$, $0.95$, $0.90$, $0.93$, $0.98$, $0.94$)
and in the predicted order up to pairs differing by less than $3\%$.  At matched $n=20$ the agreement for the long
directions improves (model$/$data $0.93$ for both $(1,7)$ and $(3,7)$ instead of $0.87$ and $0.85$ against the mixture).

\section{What fails, and what is not claimed}
\emph{A first-order pair model fails.}  The natural first attempt --- a pair of direction $v$ survives with probability
$(1-\rho)^{L}$, $\rho$ the conditional occupancy of a further cell on its line and $L$ the number of such cells, then
renormalise --- over-depletes the short directions by a factor of two ($c_{(1,1)}=0.40$ against $0.73$) and puts $(1,2)$
above $(1,1)$.  The exclusions along a long line are far from independent; the line model succeeds because it treats
each line as a unit.

\emph{The scale is imposed.}  The model's total falls $17\%$ short of \eqref{eq:total}.  After renormalisation the
fifteen measured classes sum to $17.1$ (model) against $18.6$ (data) out of the fixed total $2n-3=37$ at $n=20$; the
model's deficit in the measured directions is exactly its excess in the tail of rare, steep directions, and no second
renormalisation can move it.  Of the shortfall, at most a third can be attributed to the row and column sums that the
model ignores (a free $2n$-subset has $705.7$ non-axial pairs in expectation, the fixed-marginal model $740$);
the rest is the model's treatment of one direction at a time, whereas a solution is constrained in all directions
simultaneously.

\emph{Drift in $n$.}  The model's values drift by $3$--$8\%$ between $n=20$ and $40$ (upward for long directions),
in the same direction as the data but not by the same amount; the constancy of the short directions in the data is
reproduced only approximately ($c^{\rm model}_{(1,1)}$: $0.764$, $0.744$, $0.731$ at $n=20,30,40$).

\emph{One inversion.}  At $n=20$ the model orders $(1,6)$ above $(3,5)$ by $0.6\%$; the data order them the other way by
$4.7\%$ ($0.2492$ against $0.2609$) --- the model's largest error, $-11.5\%$, is on $(3,5)$.

\emph{Not claimed.}  No theorem about the ensemble of solutions is proved here; Theorem~\ref{thm:gf} is a statement
about a model.  No derivation of the scale, and no explanation of the growth of the long directions with $n$, is
offered.  A search of the literature beyond the classical sources (Guy and Kelly's expected-count heuristic, which
sums over primitive directions the expected number of collinear triples of a uniform $2n$-subset; Flammenkamp's
structural observations on his solution classes) has not yet been carried out; a research brief for it is on file,
and this note will be corrected if the model or the measurement is found in the literature.

\section{Reproducibility}
The measurement code, the model code with an executable statement of Theorem~\ref{thm:gf} checked against brute force
for $n\le5$, the tests written by a second agent, the two adversarial reviews (one of which found the false claim
``the data are constant in $n$'' in an earlier draft and forced the comparison at matched $n$), and the dated ledger of
predictions are public: \url{https://github.com/iwasborninbali/lemma-atelier} (lemma 005, need 008) and
\url{https://github.com/iwasborninbali/saturation}.  The database is Flammenkamp's
(\url{https://wwwhomes.uni-bielefeld.de/achim/no3in/readme.html}).

\section*{AI/tool-use disclosure}
The model, the programs and the text were produced by autonomous AI agents (Anthropic Claude) under the direction
of the author of record, who posed the question, chose what to compute, decided what to claim and takes responsibility
for the content.  Two agents worked independently --- one built the model, the other measured the database with its
own code and did not read the model before the blind test --- and a third, adversarial pass reviewed the claims
before publication.
\end{document}
